\chapter{Vergleiche und Charakterisierungen}
\label{b15guide:ch:comparison}

\section{Vollständige Fallunterscheidungen}

\begin{theorem}[Trichotomie]
\label{b15guide:thm:trichotomy}
Für $x,y\in A$ gilt genau eine der drei Aussagen
\[
 x<y,\qquad x=y,\qquad y<x.
\]
Insbesondere gilt
\[
 x\ne y\quad\Longleftrightarrow\quad(x<y\lor y<x).
\]
\end{theorem}

\begin{theorem}[Negation eines Vergleichs]
\label{b15guide:thm:negation}
Für $x,y\in A$ gelten
\[
\begin{aligned}
 x\not\leq y&\Longleftrightarrow y<x,&
 x\not<y&\Longleftrightarrow y\leq x,\\
 x\leq y&\Longleftrightarrow\neg(y<x),&
 x<y&\Longleftrightarrow\neg(y\leq x).
\end{aligned}
\]
\end{theorem}

\begin{theorem}[Einfügen eines dritten Elements]
\label{b15guide:thm:split}
Für $x,y,z\in A$ gilt
\[
 x<z\quad\Longrightarrow\quad (x<y\lor y<z).
\]
Die beiden Alternativen dürfen gleichzeitig gelten. Außerdem gilt
die negative Transitivität der strikten Ordnung:
\[
 \bigl(\neg(x<y)\land\neg(y<z)\bigr)\Rightarrow\neg(x<z).
\]
\end{theorem}

\section{Erkennen totaler Ordnungen}

\begin{theorem}[Charakterisierung durch die strikte Relation]
\label{b15guide:thm:strict-characterization}
Ist $\leq$ eine totale Ordnung, so ist die daraus definierte
Relation $<$ transitiv und erfüllt die Trichotomie.
Umgekehrt sei $\prec$ eine transitive Relation auf einer Menge
$P$, für die bei jedem Paar $x,y\in P$ genau eine der Aussagen
\[
 x\prec y,\qquad x=y,\qquad y\prec x
\]
gilt. Dann definiert
\[
 x\preceq y\quad\Longleftrightarrow\quad(x\prec y\lor x=y)
\]
eine totale Ordnung. Die beiden Übergänge liefern jeweils
die ursprüngliche Relation zurück.
\end{theorem}

\begin{definition}[Kleinste, größte und unvergleichbare Elemente]
\label{b15guide:def:extreme}
In einer geordneten Menge $(P,\preceq)$ und für $S\subseteq P$
heißt $m$ \emph{kleinstes Element} oder \emph{Minimum} von $S$,
wenn $m\in S$ und $m\preceq s$ für alle $s\in S$ gilt.
Entsprechend ist $M\in S$ ein \emph{größtes Element} oder
\emph{Maximum}, wenn $s\preceq M$ für alle $s\in S$ gilt.
Existierende Werte heißen $\min S$ und $\max S$.

Zwei Elemente heißen \emph{unvergleichbar}, wenn weder
$x\preceq y$ noch $y\preceq x$ gilt. Eine \emph{Antikette}
ist eine Teilmenge, deren verschiedene Elemente unvergleichbar sind.
\end{definition}

\begin{theorem}[Paar- und Antikettenkriterien]
\label{b15guide:thm:pair-criteria}
Für eine halbgeordnete Menge $(P,\preceq)$ sind äquivalent:
\begin{enumerate}
 \item $\preceq$ ist eine totale Ordnung.
 \item Jede Paarmenge $\{x,y\}$ mit $x,y\in P$ besitzt ein Minimum.
 \item Jede Paarmenge $\{x,y\}$ mit $x,y\in P$ besitzt ein Maximum.
 \item Jede Antikette in $P$ enthält höchstens ein Element.
\end{enumerate}
Die Paarmengen dürfen einelementig sein. Alle vier Bedingungen
gelten auch für den leeren Träger.
\end{theorem}

\begin{theorem}[Kriterium über endliche Teilmengen]
\label{b15guide:thm:finite-criterion}
Für eine halbgeordnete Menge $(P,\preceq)$ sind äquivalent:
\begin{enumerate}
 \item $\preceq$ ist total.
 \item Jede nichtleere endliche Teilmenge von $P$ besitzt ein Minimum.
 \item Jede nichtleere endliche Teilmenge von $P$ besitzt ein Maximum.
\end{enumerate}
Dabei braucht $P$ selbst nicht endlich zu sein.
\end{theorem}
